% Fragment for insertion into the book (no preamble).
% Requires only standard LaTeX; no extra packages.

\section{Two roles of classical mechanics}
\label{sec:two-roles}

Classical mechanics plays two distinct roles in physics, and it is easy to
confuse them.

In the first role it is a subject in its own right: the theory of real
macroscopic mechanical systems. As such it must accommodate everything
such systems actually do --- rigid bodies and friction, nonholonomic
constraints, bodies that change their shape, open systems, chains and
ropes. This class of systems is wide and untidy, and it has its own
``strange things'' that no appeal to a more fundamental theory can remove.

In the second role it is the base on which relativistic and quantum
theories are built. Canonical quantization replaces Poisson brackets by
commutators; relativization turns a particle into a worldline subject to
a mass-shell condition. In this role we deliberately consider only a
restricted class of systems: a finite number of degrees of freedom,
conservative forces, holonomic constraints, a regular Lagrangian and a
smooth phase space.

It would be natural to conclude that, in the second role, classical
mechanics is merely a warm-up before the ``real'' theory. This is not so.
The restriction of the class of systems is exactly what makes the problem
general: we choose the simplest systems, but the question we pose is much
broader. Which structures of the classical theory survive the passage to
a relativistic or quantum description, and which have no counterpart in
the ``parent'' theory at all? Even this narrow slice already contains
difficulties that are usually blamed on the generalizations:

\begin{itemize}
  \item \emph{Constraints and singular Lagrangians.} Dirac's
    classification of constraints and the problem of time
    (reparametrization invariance) are classical in origin.
  \item \emph{Geometric phases.} Hannay's angle is the classical
    counterpart of Berry's phase.
  \item \emph{Ambiguity of quantization.} Operator ordering, and the
    Groenewold--van Hove theorem, which shows that Poisson brackets of
    all polynomial observables cannot be consistently replaced by
    commutators, are statements about the structure of the classical
    phase space.
  \item \emph{Relativistic extension.} Radiation reaction, the
    self-energy of a point charge and Born rigidity are classical
    problems that relativity merely makes visible.
\end{itemize}

Thus the two roles must be kept apart. The first asks what classical
mechanics says about the real systems it was built for; the second asks
what a small, well-behaved part of it tells us about the relation between
theories. In the cartographic language used later in this book, the
second role fixes one particular projection of classical mechanics, while
the first shows what that projection leaves out.
